\subsection{有序子集。有序集的乘积}

设 E 是一个由序关系 \(\Gamma\) 排序的集合，其图记为 G。对于 E 的每个子集 A， \(G \cap (A \times A)\) 是 A 上的一个序；相应的序关系等价于“ \(x \leqslant y\) 并且 \(x \in A\) 并且 \(y \in A\) ”，并且我们简记为 \(x \leqslant y\) （用语上的滥用）。这样在 A 上定义的序和序关系称为由 E 上给定的序和序关系诱导的；而 E 上的序和序关系称为它们在 A 上诱导的序和序关系的扩张。每当我们把 A 视为有序集时，除非另有明确说明，我们指的是 E 上的序关系在 A 上诱导的序关系。

例子。包含关系 \(X \subset Y\) 在各种子集族上诱导的关系具有相当的重要性。以下是一些例子：

\begin{enumerate}
\def\labelenumi{(\arabic{enumi})}
\item
  设 E、F 是两个集合，并设 \(\Phi(E, F)\) 是 E 的子集到 F 的所有映射的集合。对于每个 \(f \in \Phi(E, F)\) 设 \(G_f\) 为 \(f\) 的图，它是 \(E \times F\) 的一个子集。如果我们赋予 \(\Phi(E, F)\) 序关系“ \(g\) 是 \(f\) 的延拓”（在 \(f\) 与 \(g\) 之间）（第 1 号，例 3），则 \(f \to G_f\) 是序集 \(\Phi(E, F)\) 到 \(\mathfrak{P}(E \times F)\) 的一个子集（按包含关系排序）上的同构。
\item
  对于集合 E 的每个划分 \(\varpi\) ，令 \(\tilde{\varpi}\) 为 \(\varpi\) 在 E 上定义的等价关系的图。映射 \(\varpi \to \tilde{\varpi}\) 是集合 \(\mathfrak{L}\) （E 的划分所成的集合）按关系“ \(\varpi\) 比 \(\varpi'\) 更细”（在 \(\varpi\) 与 \(\varpi'\) 之间）（第1号，例4）排序，到 \(\mathfrak{P}(E \times E)\) 的一个子集（按包含关系排序）上的同构。
\item
  设 E 是一个集合，并设 \(\Omega \subset \mathfrak{P}(\mathrm{E} \times \mathrm{E})\) 为E上的预序图集合（第2号）（或按语言习惯，E上所有预序的集合）。序关系 \(s \subset t\) （ \(s\) 与 \(t\) 之间），诱导于 \(\Omega\) （由 \(\mathfrak{P}(\mathrm{E} \times \mathrm{E})\) 上的包含关系），用以下说法来表达：“预序 \(s\) 比 \(t\) 更细”（或者说“t 比 \(s\) 更粗”）。设 \(x(s)y\) 与 \(x(t)y\) 分别表示 \((x,y) \in s\) 与 \((x,y) \in t\) 在E上的预序关系；那么说 \(s\) 比 \(t\) 更细等价于说关系 \(x(s)y\) 蕴含 \(x(t)y\) .
\end{enumerate}

设 \((\mathbf{E}_t)_{t \in I}\) 是一个集合族，并且对每个指标 \(t \in I\) 设 \(\Gamma_t\) 是 \(\mathbf{E}_t\) 上的一个序；令 \(\mathbf{G}_t \subset \mathbf{E}_t \times \mathbf{E}_t\) 是其图像，并令 \(x_t \leqslant y_t\) 表示序关系 \((x_t, y_t) \in \mathbf{G}_t\) （在 \(\mathbf{E}_t\) 上）。在乘积集合 \(\mathbf{F} = \prod_{i \in I} \mathbf{E}_t\) 上，关系

\[
(\forall \iota) ((\iota \in I) \Longrightarrow (x _ {\iota} \leqslant y _ {\iota}))
\]

是 \(x = (x_{t})\) 与 \(y = (y_{t})\) 之间的一个序关系，这一点立即可以验证。这样在F上定义的序和序关系分别称为序 \(\Gamma_{t}\) 的积和序关系 \(x_{t} \leqslant y_{t}\) 的积；这个关系记作 \(x \leqslant y\) ，而按序 \(\Gamma_{t}\) 的积排序的集合F称为有序集 \(E_{t}\) 的乘积.

立即验证可知， \(\mathbf{F}\) 上乘积序的图是乘积集合 \(\prod_{i\in I}G_i\) 在从 \(\prod_{i\in I}(E_i\times E_i)\) 到 \(\mathbf{F}\times \mathbf{F}\) 上的典范映射下的像(第二章，§5，第5号)。

有序集乘积的一个重要例子是集合 \(F^{E}\) （将集合E映入有序集F的映射的图所成的集合）。存在一个从 \(F^{E}\) 到将E映入F的映射所成的集合 \(\mathcal{F}(E, F)\) 上的典范双射，且该映射是有序集 \(F^{E}\) 到 \(\mathcal{F}(E, F)\) 上的同构，后者赋予由关系“对所有 \(x \in E\) , \(f(x) \leqslant g(x)\) ”（在从E到F的两个映射f、g之间）定义的序。该关系写作 \(f \leqslant g\) .

应当注意，在有序集 \(\mathcal{F}(\mathrm{E},\mathrm{F})\) 中，关系 \(f < g\) 意指

\pandocbounded{\includegraphics[keepaspectratio]{images/part002_2a386891d62f089ad1c0e080d8072868eea932699f46973e31691c72a41ba0df.jpg}}

``对所有 \(x \in E\) , \(f(x) \leqslant g(x)\) ，并且存在

且不

\[
f (y) <   g (y) ^ {\prime \prime}
\]

\[
" \text { 对   所有 } x \in \mathrm{E}, f (x) <   g (x)".
\]

为了避免混淆，在这种情况下最好不要使用记号 \(f < g\) 。

本小节的定义在不作改动的情况下适用于预序集，只需将全文中的“序”替换为“预序”。

